\subsection{Flujo de compresión y descompresión}

La compresión propuesta se entiende como un pipeline de etapas reversibles sobre cada bloque. Primero, el bloque en doble precisión se organiza como
un arreglo contiguo de valores. Luego puede aplicarse una transformación previa sobre la representación binaria de estos datos, por ejemplo una
variante de \emph{shuffle} o \emph{bitshuffle}, con el objetivo de exponer regularidades entre bytes o planos de bits. Finalmente, el flujo transformado
se entrega a un compresor sin pérdida de baja latencia.

El proceso inverso ocurre solo cuando el bloque es requerido por una operación. El bloque comprimido se descomprime, se invierte la transformación
previa y se obtiene nuevamente el arreglo denso de amplitudes listo para ser manipulado por la lógica del simulador. La clave del diseño es que esta
secuencia solo se ejecuta sobre los bloques activos y no sobre la totalidad del vector.

Para operaciones que afectan amplitudes pertenecientes a un mismo bloque, la actualización es enteramente local. Cuando la operación cruza bloques,
la arquitectura requiere adquirir simultáneamente los bloques involucrados, descomprimir ambos, ejecutar la actualización y decidir luego cuáles
mantener temporalmente en caché y cuáles devolver al almacenamiento comprimido. Este es el punto donde el tamaño de la caché y la granularidad de los
bloques influyen de manera más visible sobre el costo temporal.